% !TeX root = ../教材题库_全部条目_转写工作稿.tex
% Source: XJ-B1-C03-EX-P071-Y1368
% 第3章 / 3.1 函数 / 教材 p.71、72 / PDF 第78、79页
% OCR未校对：公式与文字均可能有误；未完成数学TeX转写；禁止直接发给学生
% 原题及相邻原解截图（核对用）：../题目原图/XJ-B1-C03-EX-P071-Y1368-01.webp, ../题目原图/XJ-B1-C03-EX-P071-Y1368-02.webp
% 本文件只在转写完成后作为题干使用；OCR 条目尚不是可靠题干。
\begin{bankitem}{3.1 例题 例题5}{第3章 / 3.1 函数 / 教材 p.71、72 / PDF 第78、79页}{XJ-B1-C03-EX-P071-Y1368}{OCR未校对：公式与文字均可能有误；未完成数学TeX转写；禁止直接发给学生}
\begin{ocrdraft}
[教材p.71 / PDF第78页]
\par
00作出函数3-()e[0，16]）的图象．
\par
解选择方便计算的几个数值，列表如下：
\par
1
\par
9
\par
0
\par
1
\par
4
\par
9
\par
4
\par
4
\par
4
\par
1
\par
3
\par
5
\par
0
\par
1
\par
2
\par
3
\par
4
\par
2
\par
2
\par
2
\par
[教材p.72 / PDF第79页]
\par
根据表中数据在平面直角坐标系中描点、连线，得到图3，1
\par
一2．
\par
5
\par
4
\par
3
\par
你能根据图3彐一2
\par
2
\par
估计出而，，7仃的近
\par
似值吗？
\par
0（L25]
\par
2．25
\par
4
\par
6．25
\par
9
\par
16
\par
图3．1一2
\par
函数的三种表示方法，各有长处和不足，在对具体问题的研究中需要有针对性
\par
地选择适当的表示方法，有时需要将多种方法综合运用．
\par
在数学课程中，我们主要讨论的是解析法和图象法，并且经常把这两种方法结
\par
合起来讨论．有了函数的解析式，用计算机等设备作函数图象或列函数表很方便，
\par
所以获取函数的解析式和计算方法常常是研究函数的基础性工作．
\par
当然，表示函数还有其他方法，如汽车司机要知道发动机里有多少机油，用把
\par
带刻度的尺子放进去一量便知．因为油量是油深的函数，这个函数关系体现在尺子
\par
刻度上，这种表示函数的方法称为标尺法．
\par
还如，科学研究中用到的大量函数，现在常用算法或程序来表示，而不是用列
\par
表法来表示．
\par
函数还可以用描述法来表示．例如，当工为有理数时，让D(x）一1，否则让
\par
D(x）一0，这就定义了函数D(x).它叫作狄利克雷函数．
\end{ocrdraft}
\end{bankitem}
